\section*{第 \ref{ch:b1:periodicity} 章——周期性}

\begin{solution}{exo:b1:periodicity:1}
\ce{Ca} $[\ce{Ar}]\,4\text{s}^2$：\omterm{def:b1:periodicity:table}{周期} 4，族 2，s 区。\ce{Fe}
$[\ce{Ar}]\,3\text{d}^6 4\text{s}^2$：\omterm{def:b1:periodicity:table}{周期} 4，族 8，d 区。\ce{Se}
$[\ce{Ar}]\,3\text{d}^{10} 4\text{s}^2 4\text{p}^4$：\omterm{def:b1:periodicity:table}{周期} 4，族 16，
p 区。\ce{I} $[\ce{Kr}]\,4\text{d}^{10} 5\text{s}^2 5\text{p}^5$：\omterm{def:b1:periodicity:table}{周期} 5，
族 17，p 区。
\end{solution}

\begin{solution}{exo:b1:periodicity:2}
\ce{Na} 比 \ce{Cl} 大：同\omterm{def:b1:periodicity:table}{周期}，\omterm{def:b1:periodicity:screening}{有效核电荷} $Z^{*}$ 更小。
\ce{Na} 比 \ce{Na+} 大，后者已失去 3s \omterm{def:b1:quantum-numbers:orbital}{电子层}。\ce{Cl-}
（\qty{181}{pm}）比 \ce{F-}（\qty{133}{pm}）大：多一个\omterm{def:b1:quantum-numbers:orbital}{电子层}。
\ce{Na+}（\qty{102}{pm}）比 \ce{Mg^{2+}}（\qty{72}{pm}）大：
两者电子数均为 10，核电荷数分别为 11 与 12。
% ledger: rion:Cl-VI, rion:F-VI, rion:Na+VI, rion:Mg2+VI
\end{solution}

\begin{solution}{exo:b1:periodicity:3}
$\ce{Na} (5.14) < \ce{Al} (5.99) < \ce{Mg} (7.65) < \ce{S} (10.36) <
\ce{P} (10.49) < \ce{Ar} (15.76)$。沿\omterm{def:b1:periodicity:table}{周期}的一般递增趋势被打破两次：
\ce{Al} 低于 \ce{Mg}（所移去的是第一个
3p 电子，能量高于 3s）；\ce{S} 低于 \ce{P}（所移去的
是第一个成对的 3p 电子）。
\end{solution}

\begin{solution}{exo:b1:periodicity:4}
\omterm{def:b1:periodicity:polarisability}{极化率}衡量电子云在电场作用下变形的难易程度。
\ce{I-} 比 \ce{F-} 更易极化：它更大，
且其最外层电子离原子核更远。\ce{Cl-}（含 18 个电子，
由 $Z = 17$ 束缚）的\omterm{def:b1:periodicity:polarisability}{极化率}远大于 \ce{Na+}（含 10 个电子，
由 $Z = 11$ 束缚），阴离子比阳离子更易极化，这同样成立。
\end{solution}

\begin{solution}{exo:b1:periodicity:5}
比值依次为 $18.83/5.99 = 3.1$，$28.45/18.83 = 1.5$，$119.99/28.45 = 4.2$。
跃变出现在第三个电子之后：3 个\omterm{def:b1:quantum-numbers:configuration}{价电子}，族 13，
\omterm{def:b1:periodicity:table}{周期} 3：铝。形成 \ce{Al^{3+}}，$1\text{s}^2 2\text{s}^2
2\text{p}^6$。
% ledger: ie:Al, ie2:Al, ie3:Al, ie4:Al
\end{solution}

\begin{solution}{exo:b1:periodicity:6}
加入 \ce{F-} 的电子进入较小的 2p \omterm{def:b1:quantum-numbers:orbital}{电子层}，在那里
它受到已有的 7 个\omterm{def:b1:quantum-numbers:configuration}{价电子}的强烈排斥；在氯的
较大 3p \omterm{def:b1:quantum-numbers:orbital}{电子层}中排斥较小，因此释放的能量更多。
在氮中（$2\text{p}^3$，\omorbs{u,u,u}），额外电子
必须在半充满的亚层中配对；在镁中（$3\text{s}^2$）
则须开始填入能量更高的 3p 亚层。两种情形下，
阴离子都不比原子和自由电子更稳定。
\end{solution}

\begin{solution}{exo:b1:periodicity:7}
$\chi_M(\ce{Na}) = (5.14 + 0.55)/2 = \qty{2.85}{eV}$; $\chi_M(\ce{Cl}) =
(12.97 + 3.61)/2 = \qty{8.29}{eV}$。差值很大：
\ce{Na-Cl} 键的电子对几乎完全属于氯，
于是氯化钠是离子化合物，\ce{Na+ Cl-}。
\end{solution}

\begin{solution}{exo:b1:periodicity:8}
\ce{H-F}：$\Delta\chi = 1.78$，处于离子型的边界（
实际上是极性很强的共价键），\ce{F^{\delta-}}。\ce{C-H}：0.35，
近乎非极性。\ce{Na-Cl}：2.23，离子型，\ce{Cl^{-}}。\ce{C-Cl}：0.61，
极性共价，\ce{C^{\delta+}-Cl^{\delta-}}。\ce{O-H}：1.24，极性
共价，\ce{O^{\delta-}-H^{\delta+}}。
\end{solution}

\begin{solution}{exo:b1:periodicity:9}
沿族向下，\omterm{def:b1:quantum-numbers:configuration}{价电子}处于 $n$ 更大的\omterm{def:b1:quantum-numbers:orbital}{电子层}中，离原子核更远，
并受更多\omterm{def:b1:quantum-numbers:configuration}{内层电子}\omterm{def:b1:periodicity:screening}{屏蔽}；
因此更易被移去。铷多一个\omterm{def:b1:quantum-numbers:orbital}{电子层}，\omterm{def:b1:periodicity:ionisation-energy}{第一电离能}低于
\qty{4.34}{eV}。
\end{solution}

\begin{solution}{exo:b1:periodicity:10}
锌为 $[\ce{Ar}]\,3\text{d}^{10} 4\text{s}^2$：其电子来自一个
已充满的 4s 亚层；镓为 $\dots 4\text{s}^2 4\text{p}^1$，失去其唯一的
4p 电子，能量较高，且受充满的 4s 与 3d 的\omterm{def:b1:periodicity:screening}{屏蔽}。砷为
$4\text{p}^3$，失去半充满亚层中的一个电子；硒为
$4\text{p}^4$，失去配对的电子，该电子受其配对伙伴的排斥而
能量升高，故其电离能略低。
\end{solution}

\begin{solution}{exo:b1:periodicity:11}
$\Delta E = (3.98 - 2.20)^2 = 1.78^2 = \qty{3.17}{eV} = 3.17 \times 96.5 =
\qty{306}{kJ/mol}$。然后 $D(\ce{H-F}) = \tfrac12(436 + 158) + 306 = 297 +
306 \approx \qty{603}{kJ/mol}$。较大的超出量来自
部分电荷 $\ce{H^{\delta+}}$ 与 $\ce{F^{\delta-}}$ 之间的额外吸引，
这是因\omterm{def:b1:periodicity:electronegativity}{电负性}差值较大所致。
\end{solution}

\begin{solution}{exo:b1:periodicity:12}
\omterm{def:b1:periodicity:electronegativity}{鲍林标度}由 A--B 键的键能定义：不形成键的元素
没有\omterm{def:b1:periodicity:electronegativity}{鲍林标度}值。氦、氖和氩不形成
稳定化合物。氪，尤其是氙，其价\omterm{def:b1:quantum-numbers:orbital}{电子层}更大、
更易极化，电离能也较低
（氪为 \qty{14.0}{eV}，氩为 \qty{15.76}{eV}）；它们
可被氟和氧氧化，并生成如 \ce{XeF2} 之类的化合物，
因此其键能可以测定。
% ledger: ie:Kr, ie:Ar
\end{solution}

\begin{solution}{pb:b1:periodicity:1}
\textbf{1.} 每个电子都取自比前一个离子带更多正电荷、
且半径更小的离子，后者对电子的束缚更强。
\textbf{2.} $15.04/7.65 = 1.97$；$80.14/15.04 = 5.33$；$109.27/80.14 = 1.36$。
\textbf{3.} 跃变出现在移去两个电子之后：X 有 2 个\omterm{def:b1:quantum-numbers:configuration}{价电子}，
族 2；在第 3 \omterm{def:b1:periodicity:table}{周期}中它是镁。
\textbf{4.} \ce{Mg^{2+}}，$1\text{s}^2 2\text{s}^2 2\text{p}^6$，
即氖的\omterm{def:b1:quantum-numbers:configuration}{电子排布}。
\textbf{5.} 金属：属于族 2，位于\omterm{def:b1:periodicity:table}{周期}表左侧，
\omterm{def:b1:periodicity:ionisation-energy}{第一电离能}较低。
\textbf{6.} 第三个电子来自 2p 内层，其 $n$ 更小，
离原子核近得多，几乎不被\omterm{def:b1:periodicity:screening}{屏蔽}，且是从
带两个正电荷的离子中移去的。
% ledger: ie:Mg, ie2:Mg, ie3:Mg, ie4:Mg
\textbf{7.} 均有 10 个电子：$1\text{s}^2 2\text{s}^2 2\text{p}^6$。
\textbf{8.} 相同的 10 个电子分别受核电荷 8、9、11、
12、13 的束缚：随着 $Z$ 增大，电子云收缩。
\textbf{9.} \ce{O^{2-}}：半径最大，对其
10 个电子而言核电荷最低，且为阴离子。
\textbf{10.} $140/54 = 2.6$。
\textbf{11.} $198.8/2 = \qty{99.4}{pm}$；阴离子约为
\omterm{def:b1:periodicity:radius}{共价半径}的 $181/99.4 = 1.8$ 倍。
% ledger: re:Cl2, rion:Cl-VI
\textbf{12.} $\ce{Mg^{2+}} < \ce{Na+} < \ce{Cl-}$：两种阳离子
互为等电子体（核电荷分别为 12 与 11），而 \ce{Cl-} 多一个\omterm{def:b1:quantum-numbers:orbital}{电子层}。
\textbf{13.} $\chi_M$ = 10.41 (\ce{F}), 8.29 (\ce{Cl}), 7.59 (\ce{Br}),
7.53（\ce{O}）、6.26（\ce{C}），单位为 eV。
\textbf{14.} 马利肯：$\ce{F} > \ce{Cl} > \ce{Br} > \ce{O} > \ce{C}$；
鲍林：$\ce{F} > \ce{O} > \ce{Cl} > \ce{Br} > \ce{C}$。氧
由第四位升至第二位。
\textbf{15.} $3.98/10.41 = 0.382$, $3.16/8.29 = 0.381$, $2.96/7.59 =
0.390$：近似为常数，约 0.38；对卤素而言，两种标度
成正比。
\textbf{16.} 氧的\omterm{def:b1:periodicity:electron-affinity}{电子亲和能}较小（\qty{1.44}{eV}），
因为所加电子必须在紧凑的 2p \omterm{def:b1:quantum-numbers:orbital}{电子层}中配对；
孤立原子的马利肯值低估了氧在键中施加的
拉力，而由键构建的\omterm{def:b1:periodicity:electronegativity}{鲍林标度}则直接测量这种拉力。
\textbf{17.} \ce{HCl} 中的氯；\ce{ClF} 中的氟。
\textbf{18.} $\Delta\chi = 0.96$：极性共价键。
\textbf{19.} $D(\ce{H-H}) = 2 \times 218.0 = \qty{436.0}{kJ/mol}$。
\textbf{20.} $D(\ce{Cl-Cl}) = 2 \times 121.3 = \qty{242.6}{kJ/mol}$。
\textbf{21.} $D(\ce{H-Cl}) = 218.0 + 121.3 - (-92.3) = \qty{431.6}{kJ/mol}$。
\textbf{22.} $\Delta E = 431.6 - \tfrac12(436.0 + 242.6) = 431.6 - 339.3 =
\qty{92.3}{kJ/mol}$。用符号表示，$D(\ce{H-Cl}) = \Delta_fH(\ce{H}) +
\Delta_fH(\ce{Cl}) - \Delta_fH(\ce{HCl})$，而对称键的
平均值为 $\Delta_fH(\ce{H}) + \Delta_fH(\ce{Cl})$，因此 $\Delta E =
-\Delta_fH(\ce{HCl})$。
\textbf{23.} $92.3/96.485 = \qty{0.957}{eV}$。
\textbf{24.} $|\chi(\ce{Cl}) - \chi(\ce{H})| = \sqrt{0.957} =
\textbf{0.98}$，对比表中的 $3.16 - 2.20 = 0.96$：鲍林的
运算采用现代数据，所得差值与表中数值
相差在 0.02 之内。
\end{solution}
